\chapter{Actuation: from PWM to force}
\label{ch:actuation}

The thrust $T$ in \eqref{eq:surge} is never measured. It is \emph{reconstructed} from the
logged servo PWM and the battery voltage. Any error in this reconstruction goes straight
into the identified coefficients (\cref{sec:identifiability}), so this chapter matters as much
as the hydrodynamics.

\section{The signal chain}

\begin{enumerate}
  \item A command script publishes \code{Int16MultiArray} $[x, y, z, r, s, t]$ on
        \code{/pixhawk/manual_control}, with surge $x\in[-1000, 1000]$, capped at $\pm500$.
  \item The bridge forwards it as a MAVLink \code{MANUAL_CONTROL} message. ArduSub in MANUAL mode passes it
        through the motor mixer, without any closed-loop controller, to produce one PWM per motor.
  \item The ESC drives the motor, the propeller produces thrust, and the PWM actually sent is
        echoed on \code{/pixhawk/servo_output_raw} (6 channels, \unit{\micro\second}, no header).
\end{enumerate}
The model starts at step 3. It uses the PWM that was \emph{sent} to the ESC, not the command,
so the mixer and any ArduSub limits are already accounted for.

\section{Thruster physics and the T500 table}

A propeller's static thrust scales with the square of its rotation rate,
$T = \rho K_T n^2 D^4$, and the rotation rate an ESC can reach rises with the supply voltage.
The same PWM therefore produces more thrust from a full battery than from a depleted one.
Blue Robotics publishes measured T500 thrust against PWM at several voltages. The repository
keeps four columns (\SIlist{12;14;16;18}{\volt}) in \file{thruster_curves/t500_force_vs_voltage.csv},
plotted in \cref{fig:t500}.

\begin{figure}[htbp]
\centering
\begin{subfigure}{0.48\textwidth}
\centering
\begin{tikzpicture}
\begin{axis}[notes,xlabel={PWM / \unit{\micro\second}},ylabel={Thrust / \unit{\newton}},xmin=1100,xmax=1900,
  xtick={1100,1300,1500,1700,1900}]
\addplot[blue!35] table[x=pwm,y=f12]{data/t500.dat};
\addplot[blue!60] table[x=pwm,y=f14]{data/t500.dat};
\addplot[blue!85!black] table[x=pwm,y=f16]{data/t500.dat};
\addplot[black] table[x=pwm,y=f18]{data/t500.dat};
\legend{\qty{12}{\volt},\qty{14}{\volt},\qty{16}{\volt},\qty{18}{\volt}}
\end{axis}
\end{tikzpicture}
\caption{Thrust against PWM}
\end{subfigure}\hfill
\begin{subfigure}{0.48\textwidth}
\centering
\begin{tikzpicture}
\begin{axis}[notes,xlabel={Battery voltage / \unit{\volt}},ylabel={Thrust / \unit{\newton}},xmin=11,xmax=19,ymin=0]
\addplot[blue!40] table[x=v,y=f1700]{data/t500_volt.dat};
\addplot[blue!75!black] table[x=v,y=f1800]{data/t500_volt.dat};
\addplot[black] table[x=v,y=f1900]{data/t500_volt.dat};
\addplot[gray,dashed,forget plot] coordinates {(14,0) (14,120)};
\addplot[gray,dashed,forget plot] coordinates {(16.8,0) (16.8,120)};
\legend{\qty{1700}{\micro\second},\qty{1800}{\micro\second},\qty{1900}{\micro\second}}
\end{axis}
\end{tikzpicture}
\caption{Thrust against voltage (\code{ThrusterMap})}
\end{subfigure}
\caption{Blue Robotics T500 thrust from the manufacturer table: against PWM at the four
tabulated voltages, and against battery voltage at three PWM values as evaluated by
\code{ThrusterMap.force_n}. Dashed lines mark \qty{14.0}{\volt} (BMS warning threshold) and
\qty{16.8}{\volt} (full 4S pack). Outside \qtyrange{12}{18}{\volt} the map is clamped.}
\label{fig:t500}
\end{figure}
\FloatBarrier

Four features of the table matter for the model:
\begin{itemize}
  \item \textbf{Dead band.} Thrust is zero in a band around the \qty{1500}{\micro\second} neutral (about
        \qtyrange{1460}{1540}{\micro\second}; the exact edges differ slightly between voltages).
        Small commands produce no force at all.
  \item \textbf{Voltage sensitivity.} At \qty{1900}{\micro\second} thrust rises from
        \qty{58.1}{\newton} at \qty{12}{\volt} to \qty{104.5}{\newton} at \qty{18}{\volt}, roughly
        $T\propto V^{1.45}$. Over a normal dive (\qty{16.8}{\volt} down to \qty{14}{\volt})
        full-scale thrust falls by about \qty{25}{\percent}. A model evaluated at a fixed voltage
        would absorb that drift into its drag coefficients.
  \item \textbf{Forward/reverse asymmetry.} At \qty{16}{\volt}, \qty{1900}{\micro\second} gives
        \qty{89.8}{\newton} but \qty{1100}{\micro\second} gives only \qty{-59.8}{\newton}, because
        propellers are optimised for one direction.
  \item \textbf{Units.} The table is in kilogram-force. The code multiplies by
        $g = \qty{9.80665}{\meter\per\second\squared}$ (\code{KGF_TO_N}).
\end{itemize}

\section{Bilinear interpolation}

\code{ThrusterMap.force_kgf_at(p, V)} first interpolates linearly in PWM inside every voltage
column, then linearly in voltage between the two bracketing columns. With $V_j \le V \le V_{j+1}$:
\begin{equation}
  F(p, V) = (1-\lambda)\,F_j(p) + \lambda\,F_{j+1}(p), \qquad
  \lambda = \frac{V - V_j}{V_{j+1} - V_j},
  \label{eq:bilinear}
\end{equation}
where $F_j(p)$ is the PWM interpolation in column $j$. For example
$F(\qty{1700}{\micro\second}, \qty{15.2}{\volt}) = \qty{2.998}{kgf} = \qty{29.4}{\newton}$.
Inputs outside the table are \emph{clamped} to the nearest edge and never extrapolated, and
\code{clamped()} reports when that happened. \file{propagate_model.py} prints a warning when
voltage samples fall outside the \qtyrange{12}{18}{\volt} range.

\section{Thrust direction: mirroring, not negating}

If a thruster turns out to be mounted backwards (\code{thrust.sign: -1}), the code does
\emph{not} compute $-F(p)$. It mirrors the PWM about neutral instead:
\begin{equation}
  F_\mathrm{eff} = F\bigl(p_0 + s\,(p - p_0)\bigr), \qquad s = \pm1 .
  \label{eq:mirror}
\end{equation}
With an asymmetric curve the two are not equivalent. A reversed propeller pushed by a
``forward'' PWM produces the \emph{reverse} branch of the curve, not the negated forward
branch. \code{_mirror()} in \file{actuation.py} implements \eqref{eq:mirror}. The same
mirroring is applied to the baseline power law in \file{propagate_model.py}.

\section{Battery voltage}

The BMS node publishes the pack voltage about once per second as a \code{DiagnosticStatus}
named \code{"Battery: Voltage"} on \code{/diagnostics}, with the value as text
(\code{"14.70V"}). \code{read_voltage()} parses it (\code{parse_bms_voltage}) and
\code{VoltageTrack.at()} holds it constant between samples (zero-order hold). Before the
first sample the first value is used. That is harmless because pack voltage changes over
minutes. If a bag has no BMS data, \code{thrust.voltage_fallback_v} or the \code{--voltage}
flag must provide a fixed voltage. Otherwise the scripts refuse to run.

\section{Unknown is NaN, never zero}

\code{ThrustSource} returns \code{NaN} for every servo channel that has no curve yet. A
missing curve can therefore never be mistaken for zero thrust, which would silently bias a
fit. The same policy applies to stale PWM: if the last servo sample is older than
\code{pwm_stale_s} (three times the median servo interval, at least \qty{0.5}{\second}),
the thrust over that span is unknown. The propagation then restarts from the EKF instead
of integrating with a guessed force.

\section{Allocation: from thruster forces to body forces}

Thruster $i$ at position $\vect{r}_i$ (relative to the centre of mass) pushing along unit
direction $\vect{d}_i$ with force $F_i$ contributes
\begin{equation}
  \vect{\tau}_i = \begin{bmatrix} \vect{d}_i \\ \vect{r}_i\times\vect{d}_i \end{bmatrix} F_i,
  \qquad
  \vect{\tau} = \sum_i \vect{\tau}_i = \vect{B}\,\vect{F},
  \label{eq:allocation}
\end{equation}
where the columns of the allocation matrix $\vect{B}$ are the per-thruster vectors. The
\code{allocation:} section of \file{model.yaml} stores one row of $\vect{B}$ per axis as
\code{(channel, lever_m)} pairs (\cref{tab:allocation}). For translation the ``lever'' is the
direction cosine, for rotation it is the moment arm in metres.

\begin{table}[htbp]
\centering
\caption{Allocation rows stored in \file{model.yaml} (FLU, about the centre of mass).}
\label{tab:allocation}
\begin{tabular}{@{}l l S[table-format=-1.4] l@{}}
\toprule
Axis & Channel & {Lever} & Note \\
\midrule
surge & 0 (\code{MOT_1}) & 1.0 & only thruster with forward authority \\
pitch & 2 (\code{MOT_3}) & 0.5508 & rear vertical pair \\
pitch & 3 (\code{MOT_4}) & 0.5508 & rear vertical pair \\
pitch & 4 (\code{MOT_5}) & -0.6358 & front vertical \\
pitch & 0 (\code{MOT_1}) & 0.0300 & surge thruster \qty{3}{\centi\meter} above the CoM \\
\bottomrule
\end{tabular}
\end{table}
\FloatBarrier

Only the surge row is used today, since surge is the only modelled DOF. The other channels'
curves, signs and the remaining axes' rows are still to be added, and the lever arms will
change when the Pixhawk mixer is updated for the new geometry.

\section{Thrust bench (planned)}
\label{sec:thrust-bench}

\begin{revised}[title=Added 28 October 2026,breakable]
The vertical and lateral thrusters (DD2838) have no published curve, and the curve in
\file{thruster_curves/PLACEHOLDER_dd2838_force_vs_voltage.csv} is a scaled T500 table. A bench
is planned to measure the static thrust (bollard pull) $F(p, V)$ of one representative unit.
The concept document is \emph{Thrust-Messstand: Konzept und Messablauf} (28 October 2026). Its
main points are summarised here, together with corrections from its review.

\paragraph{Set-up.} The thruster is mounted underwater on a lever or pendulum, and the force is
led out of the water to a dry S-type load cell (\qty{10}{\kilogram} rating, about
\qty{98}{\newton}), connected at both ends through rod ends so that only axial force reaches it.
The expected maximum of \qty{3.1}{kgf} (about \qty{30}{\newton}) uses about \qty{31}{\percent} of
the rating. The cell is read by an Adafruit NAU7802 (24-bit ADC, I\textsuperscript{2}C) on the
Jetson Orin Nano. The PWM is set from the same computer through a Pixhawk running ArduSub,
on an output with \code{SERVOn_FUNCTION = 0} so that \code{MAV_CMD_DO_SET_SERVO} drives it
directly. A fixed-voltage power supply feeds the ESC, and the voltage is measured at the ESC
and logged, because cable losses under load reach a few hundred \si{\milli\volt}.

\paragraph{Procedure.} One sweep is run per supply voltage, with all raw samples logged with a
time stamp; averaging and the removal of the settling phase are done afterwards, so that the
window can be changed later. \Cref{tab:bench-procedure} lists the parameters.

\paragraph{Analysis.} The ADC counts are converted to newtons with a calibration factor from
reference masses, applied through the lever, so that the lever ratio is part of the calibration.
The first \qty{5}{\second} of every step are discarded (the settling time is estimated at
hundreds of milliseconds to a few seconds and is checked by comparing windows that skip 3, 5 and
\qty{8}{\second}). Further checks: the tare before and after each sweep, up against down
(hysteresis, heating), repeated sweeps, the logged voltage against the set value, and the logged
\code{SERVO_OUTPUT_RAW} against the commanded PWM. The result is written in the format of
\file{t500_force_vs_voltage.csv} (\code{pwm_us}, \code{force_kgf_<V>V}, one common PWM grid for
all voltages), so that \code{ThrusterMap} reads it without changes.

\paragraph{Corrections from the review.}
\begin{itemize}
  \item \textbf{Excitation voltage.} The concept sets the bridge excitation to \qty{4.0}{\volt}.
        The NAU7802's internal regulator cannot produce more than its own supply minus a dropout,
        so with the board supplied from \qty{3.3}{\volt} about \qty{3.0}{\volt} is the realistic
        maximum. At \qty{3.0}{\volt} and \qty{2}{\milli\volt\per\volt} the full-scale signal is
        \qty{6}{\milli\volt} (\qty{1.8}{\milli\volt} at \qty{30}{\newton}), which is still ample.
  \item \textbf{Reverse thrust is required, not optional.} The vertical and lateral thrusters run
        in both directions on the vehicle, and the forward and reverse branches differ
        (\cref{fig:t500}).
  \item \textbf{Finer steps near neutral.} The vehicle spends most of its time near
        \qty{1500}{\micro\second}, where the dead band ends (often \qtyrange{25}{50}{\micro\second}
        either side). Steps of \qty{10}{\micro\second} within about \qty{\pm 100}{\micro\second}
        of neutral resolve where the thrust starts.
  \item \textbf{Voltages.} The sweeps should bracket the operating range of the 4S pack
        (about \qtyrange{13}{16.8}{\volt}); the interpolation clamps outside the measured range.
        The concept's scaling $F \propto V^2$ is an upper estimate: the T500 data follow about
        $V^{1.5}$.
  \item \textbf{The same ESC as on the vehicle.} The ESC sets the dead band and the mapping from
        PWM to speed, so the curve is only valid for the ESC model used on the bench.
  \item \textbf{Measured in water (decided).} The concept left water or air open; the bench is
        run in water. In air the propeller is unloaded and the motor is not cooled. In a small
        tank the jet recirculates and lowers the thrust over time; the up/down and repeat
        comparisons reveal it.
\end{itemize}
\end{revised}

\begin{table}[htbp]
\centering
\caption{Sweep parameters of the planned thrust bench, from the concept document with the
corrections of \cref{sec:thrust-bench}.}
\label{tab:bench-procedure}
\begin{tabular}{ll}
\toprule
Parameter & Value \\
\midrule
PWM steps & \qty{20}{\micro\second}; \qty{10}{\micro\second} within \qty{\pm 100}{\micro\second} of neutral \\
Range & \qtyrange{1500}{1900}{\micro\second} and \qtyrange{1500}{1100}{\micro\second}, each up and down \\
Hold per step & \qty{15}{\second}, of which the first \qty{5}{\second} are discarded \\
Ramp between steps & \qty{10}{\micro\second} per \qty{0.1}{\second} \\
Tare & \qty{10}{\second} at \qty{1500}{\micro\second} before and after each voltage \\
Pause between voltages & about \qty{60}{\second} \\
Voltages & bracketing \qtyrange{13}{16.8}{\volt}, measured at the ESC \\
\bottomrule
\end{tabular}
\end{table}
\FloatBarrier

\section{The deployed baseline}

\file{thruster_velocity_estimator.py} estimates speed as
\begin{equation}
  u_\mathrm{pl} = 0.780\,\mathrm{sign}(c)\,|c|^{0.964}, \qquad c = \frac{p - 1500}{500}.
  \label{eq:powerlaw}
\end{equation}
This is a steady-state map fitted directly from PWM to speed. Compared with
\eqref{eq:surge} it implicitly assumes (i) instantaneous response (no inertia), (ii) a fixed
voltage, (iii) no dead band, and (iv) symmetric forward/reverse behaviour. It is kept in
\code{dynamics.powerlaw_speed} solely as a comparison: a dynamic model that cannot beat it
has not earned its complexity.

\begin{questions}
\begin{enumerate}
  \item Why would identifying drag from a fixed-voltage thrust curve make the fitted $\Xuu$
        drift between the start and end of a battery?
  \item A thruster is found to be reversed. Why is $F(p_0 - (p - p_0))$ more correct than
        $-F(p)$?
  \item With the levers of \cref{tab:allocation}, what pitch moment does \qty{20}{\newton} of
        surge thrust produce, and in which direction does the nose move?
\end{enumerate}
\end{questions}
